\subsection{2.  Team Dynamics, Membership Change, and Coordination}

A growing body of organizational behavior research addresses how team composition changes affect team functioning and performance. Morgeson, Mitchell, and Liu [16] introduced Event System Theory (EST), which posits that organizational events are characterized by three properties---novelty, disruption, and criticality---that determine their impact on organizational entities. Events that are higher on these dimensions are more likely to produce changes in behavior, features, or subsequent events.

Summers, Humphrey, and Ferris [20] applied an event-based lens to team member change and found that membership changes caused high levels of flux in coordination when the change involved a strategically core role or when information transfer during the transition was low. Trainer et al. [22] reconceptualized each team membership change as a discrete team-level event with varying degrees of novelty, disruptiveness, and criticality, finding that newcomer entry can substantially impact team functioning beyond its impact on the newcomer themselves.

Research on Transactive Memory Systems (TMS)---collective memory systems in which team members specialize in different knowledge domains while maintaining shared awareness of who knows what [23, 15, 18]---provides a cognitive mechanism for understanding how membership changes disrupt coordination. When a new member enters, the team's TMS becomes inaccurate: existing members' calibrations of who knows what and how to coordinate are invalidated [15, 23]. The newcomer socialization literature further documents that support for newcomers declines within the first 90 days [6], and that structured onboarding significantly improves adjustment outcomes [2].

These streams provide the theoretical vocabulary for Section 8.3: how a new agent introduced to an existing mesh can degrade peer performance even when performing its own task well.

